% TOP_PROBLEM: {"id": 3100053, "title": "Show that there exists a B so that, for every n > 0, there exists a k relatively prime to n whose continued fraction has", "category_id": 1, "category": {"id": 1, "name": "number_theory", "display_name": "Number Theory", "description": "Properties of integers, prime numbers, Diophantine equations.", "slug": "number-theory", "order_index": 1, "created_at": "2026-07-31T15:26:25.670Z"}, "status": "open", "problem_number": "AMR-030-0053", "set_id": 13, "statusQualification": "Makes explicit the intended fraction k/n, every-prefix averaging, and the canonical rational continued-fraction convention; n=1 is vacuous."}
% FIRST_POSTED: 2026-10-02
% ENTRY_KIND: statement_only
% UnsolvedMath ID 3100053; source catalogue code AMR-030-0053
% !TeX program = lualatex
% Definition and sources only; this file is not an LLM solution attempt.
\documentclass[11pt,a4paper]{article}
\usepackage[margin=25mm]{geometry}
\usepackage{fontspec}
\setmainfont{lmroman10-regular.otf}[BoldFont=lmroman10-bold.otf,
 ItalicFont=lmroman10-italic.otf,BoldItalicFont=lmroman10-bolditalic.otf]
\usepackage{amsmath,amssymb,mathtools}
\usepackage{unicode-math}
\setmathfont{latinmodern-math.otf}
\AtBeginDocument{\let\setminus\smallsetminus}
\usepackage{xurl,hyperref}
\hypersetup{colorlinks=true,urlcolor=blue}
\setcounter{secnumdepth}{0}
\setlength{\parindent}{0pt}
\setlength{\parskip}{5pt}
\setlength{\emergencystretch}{3em}
\begin{document}
\section*{Show that there exists a B so that, for every n > 0, there exists a k relatively prime to n whose continued fraction has}\label{3100053}
{\small UnsolvedMath ID 3100053 · AMR-030-0053 · Number Theory · AMR Open Problem Lists}
\subsection{Definitions and mathematical statement}
For an integer $n\ge2$, choose $1\le k<n$ with $\gcd(k,n)=1$
and write its regular continued fraction in the canonical form
\[
       \frac{k}{n}=[0;a_1,\ldots,a_r],
       \qquad a_i\in\mathbb Z_{>0},\quad a_r\ge2.
\]
Here the notation means the finite nested fraction
$1/(a_1+1/(a_2+\cdots+1/a_r))$. Say that its partial quotients
are bounded in prefix average by $B$ if
\[
                    \sum_{i=1}^{j}a_i\le Bj
                         \qquad(1\le j\le r).
\]
Does an absolute constant $B$ exist such that every denominator
$n\ge2$ admits a coprime numerator satisfying all these inequalities?
In particular, is $B=3$ sufficient?

The numerator may depend on $n$, but the constant may not. The
condition concerns every initial segment, not only the average
over the entire continued fraction; large early quotients cannot
be compensated by later small ones. It is also weaker than
requiring each quotient separately to be at most $B$.
The source's reference to a continued fraction of a multiplier
$k$ is interpreted as the rational number $k/n$. The trivial
denominator $n=1$ is omitted, and the terminal-quotient convention
removes the two-expansion ambiguity for rational numbers.
\subsection{Short English statement}
Does every denominator admit a coprime numerator whose continued-fraction partial quotients have a uniformly bounded average on every prefix, perhaps with bound3?
\subsection{Sources}
\begin{itemize}
\item[S1] ULAM.AI and the UnsolvedMath Contributors, supplied problems.json, record 3100053 (AMR-030-0053), AMR Open Problem Lists. Catalogue identity and source transcription; CC BY 4.0.
\url{https://huggingface.co/datasets/ulamai/UnsolvedMath}.
\item[S2] Cooper - Combinatorial Problems I Like (2020 snapshot), Combinatorial Number Theory, source-order bullet 53. Original source indexed by AMR-030-0053.
\url{https://people.math.sc.edu/cooper/combprob.html}.
\end{itemize}
\end{document}
